\chapter{The Number Splits Into Faces}

\section{One corridor, four faces}

The first two parts of this book built a single thing and defended it: a residue, bracketed, floored, and repeatable
--- an honest number that reruns identically and that anyone can check. This part opens by doing something that looks,
for a moment, like taking it apart. The one number, it turns out, has more than one honest reading. Decoded in
different frames, the same residue shows different \emph{faces} --- and the faces are not rival answers competing to be
the true one. They are projections of a single object onto a basis, each reading a part of it the others cannot see.
The surprise of the part is that this splitting, far from weakening the number, is exactly where its honesty gets
proved.

Start with the structure, which is a basis. The construction does not read the residue as a bare scalar --- a lone
quantity on a line. It reads it as an element of a \emph{basis}: a decomposition into components, each a coordinate
along one direction. The type that carries this has three constructors, and the one to hold onto is the first: a
\emph{null space} --- the shared origin from which every component is measured, the common zero all the faces open off.
Call it the corridor. Every face is a reading taken down that one corridor, from that one origin, along a different
direction. To decode the residue is to express it in this basis; each face is one coordinate of the same object, read
from the same origin.

Now the property that makes the faces worth having, and it is the heart of the section. The basis carries an
\emph{orthogonality} relation --- a precise test for when two directions are decorrelated, each blind to what the other
reads. The relation is exact. The origin, the zero of the space, is orthogonal to everything --- it shares no direction
with any reading. Two base directions are orthogonal when they sit at different positions; two scaled components are
orthogonal when their leading edges differ or their underlying facts clash. In every case orthogonality means the same
thing: the two readings do not overlap, so neither can be inferred from the other. This is the discipline the whole
four-book project is built on, and it has a name --- the \emph{null basis}. Its point is to choose the readings so that
they \emph{cannot see each other}: decorrelated on purpose, so that when two of them agree, the agreement carries
genuine information and is not one reading echoing a bias it shares with the next.

Here is why that matters, and it is the reason this book has siblings. The instrument is read four ways --- this
volume, in the language of type theory and construction; and three others, in physics, in computation, and in a walk
through the code itself. Those four are not four summaries of one argument. They are \emph{four orthogonal decodings of
the one residue}, chosen deliberately to share no vocabulary --- so that when all four land on the same reading, the
coincidence cannot have been arranged. A crank can tune a single face to any number he wishes; a wished-for value can
always be made to fall out of one decoding if you build the decoding around it. What a crank cannot do is tune four
decorrelated faces to agree by accident, because to fix the number on one is to say nothing about where the others
land. So ``the number splits into faces'' is not a weakening of the claim. It is the machine's strongest form of
evidence --- agreement across independent readings, consilience --- and it is the anti-crank argument the whole
four-book project exists to make. (That there are exactly \emph{four} is the expository shape of the thing, the count
of the gauges; what the code carries is the deeper fact beneath it --- a residue decoded into an orthogonal basis,
readings that do not see each other. The number four is the book's, not the machine's.)

In the two verbs the faces have a clean description. Each face is a \emph{tange} --- a selection of one component of
the residue, picked out by the characteristic that its basis direction reads. The four faces together are a
\emph{funge} --- the readings bagged as ``the same object, seen four ways.'' And the null basis is precisely the
discipline of choosing those tanges \emph{orthogonal}, so that the funge which gathers them is honest: four genuinely
different selections, along directions that cannot see one another, that nonetheless bag to a single number. A
dishonest instrument would tange the same component four times and call the four agreements a confirmation; this one
tanges four decorrelated components, and their agreement is real.

And this is the argument the ending rests on. When the book finally hands back its bracket --- the interval floored at
$\alpha$, held open --- the reason to trust it is not that this volume's decoding is especially careful. It is that
four decodings, in four disjoint vocabularies, each blind to the others, close on the same interval. The jar is
believable because four faces that cannot see each other agree on it. The rest of this part follows the faces down:
where they cut, how the same residue reads covariantly on one band and contravariantly on the other, and why one
object read many ways is stronger evidence than one object asserted once. The number did not break when it split. It
got proved.

\section{The partition point}

The last section said the faces are orthogonal --- decorrelated, each reading what the others cannot. This section
finds \emph{where} they separate, and the answer is sharper than ``four different directions.'' At bottom the split is
not four-way at all. It is a single binary test, applied to every reading, and the whole four-fold structure hangs off
it. The test is this: \emph{do these two facts agree?} When they agree, the reading runs one way; when they disagree,
it runs the inverted way. That one fork is the partition point --- and it turns out to be exactly the seam between the
book's two working verbs.

Look at the order the basis carries, at the case where it compares two full readings. The comparison forks, and it
forks on a single condition: whether the two readings' underlying facts are equal. If the facts match, the order takes
one form; if they differ, it takes another. Nothing else in the case decides the routing --- not the magnitudes, not
the directions on their own. The fact-parity test alone --- \emph{agree} or \emph{disagree} --- sends the comparison
down one branch or the other. That disjunction is the partition point: every reading the basis makes passes through
this one hinge, and which side it comes out on is fixed by whether its facts match.

Now the two branches, because their difference is the whole point. On the agree-branch, the products scale
\emph{covariantly}: both compare in the same direction, the reading running with its terms. On the disagree-branch,
they scale \emph{contravariantly}: the comparison inverts, the order on the terms flipped, so that what was below is
now above. Same two readings, opposite verdict, and the difference is nothing but whether their facts agreed. Those two
branches are the two bands the book has been building toward. The covariant band is the reading that runs \emph{with}
the world --- the count, the parity the object itself carries. The contravariant band is the reading that runs
\emph{against} it --- ours, the labels we lay over the object, inverted relative to what it counts. The partition point
is where those two diverge: one hinge, and past it, the world's reading and our reading fork apart.

And here the two verbs arrive at their master appearance. The whole of this fork, in the book's working language, is
the tange/funge distinction made into a single test. Fact-agreement \emph{funges}: it counts \emph{with}, gathering by
the parity the world supplies --- the covariant band. Fact-disagreement \emph{tanges}: it selects \emph{against}, by
the very characteristic on which the facts differ --- the contravariant band, with its boundaries inverted. Every
earlier hint of this --- the covariant and contravariant arms the order rules carried, the double-routing the sampled
residue performed --- lands here as the master cut. The two verbs are not two moods the book adopts; they are the two
branches of a single disjunction the machine computes, and the seam between them is exactly \emph{do the facts agree}.
That is why the four faces of the last section are not four independent things. They are combinations of these two
bands --- readings sitting on the covariant side, on the contravariant side, or reading the fork itself. The partition
point is the one hinge from which all the faces open.

One thing this fork carries, the book marks and defers. The two bands --- the world's parity and our inverted labels,
the covariant and the contravariant --- are, read physically, a split the payoff of this volume turns on: matter on one
side, its mirror on the other, and a \emph{sign} that flips as you cross the fork. But the sign itself is not this
section's to give. It belongs to a later chapter --- the one that rotates the corridor onto its imaginary axis and
reads the flip as a quarter-turn --- and to the physical volume, which reads the fork as an interaction and names what
the inverted branch produces. Here the book plants only the fork: a single fact-parity test, forking covariant from
contravariant, funge from tange. Given a sign, that fork becomes the four-fold the payoff needs; without one, it is
exactly what it looks like --- the one binary hinge every face of the number is routed through, the seam where the
world's count and our labels part ways.

\section{Same residue, different reading}

Two sections have opened the number into faces and found the seam they cut along. This closing section lands the point
the whole chapter was built to make, and it is the point the entire four-book project stands on. It is \emph{one}
residue. The faces are not four residues that happen, remarkably, to agree. They are a single object, read four ways
--- and that single fact is what turns their agreement from a coincidence into evidence.

Look at what the chapter's relations actually take as their argument. The residue, decoded into the basis of the first
section, is a single object --- one inductive value, built from an origin and its products. And every relation in this
chapter reads \emph{that} object: the order on the basis, the orthogonality test, the fact-parity fork --- each of them
is a function of the same one thing. There are not four objects here quietly kept in step. There is one object and
several questions asked of it. That is the structural fact beneath everything the chapter has said: the faces multiply,
the residue does not.

Now watch the readings vary while the object holds still. The same residue-object is read by its own order --- which of
two stands below the other. It is read by the orthogonality relation of the first section --- is this direction blind
to that one, or do they overlap? It is read by the fact-parity fork of the last section --- does this comparison run on
the covariant band or the contravariant one? Three relations, and each takes the same object as its argument; what
differs between them is not the thing read but the reading. ``Same residue, different reading'' is not a figure of
speech. It is the literal shape of the code: the object is fixed, and the relation brought to bear on it varies. A face
is not a piece of the residue broken off; it is the whole residue, seen through one relation.

And here the chapter closes the argument the first section opened. Because there is exactly one object underneath, the
orthogonal faces reading it are not free to say whatever they like --- they are projections of the same thing, and
projections of one object are constrained to be consistent with each other. So when the four gauges, decorrelated by
construction so that none can see what the others read, are each turned on the one residue and each return the same
bracket, that agreement is not luck and it is not collusion. It reflects the object. This is the whole force of the
null basis, and it is worth stating plainly: a single gauge, left alone, can be tuned to wish. Four faces of one object
cannot wish in concert --- because they are not four wishes, they are one thing, refracted, and a refraction cannot
disagree with the object it refracts. Many readings, one residue, agreement: therefore the reading is real.

In the two verbs the whole chapter is a round-trip. The first section \emph{tanged} the one residue into faces ---
selecting components, one per direction. The second cut those faces at fact-parity, sorting them onto the two bands.
And this section \emph{funges} them back: the several readings bag, at the end, to a single object. What makes that
closing funge honest is exactly what the chapter has been proving --- that there is one thing to bag. A dishonest
instrument would tange four faces out of nothing and funge four wishes into a fake agreement; this one tanges four
faces of a real object and funges them back onto it, and the fidelity of the round-trip is the evidence. Tange into
orthogonal faces, funge back to the one residue: when the bag closes on the object it started from, the reading has
proved itself.

So the seventh chapter ends where it aimed to. The number did not fracture into rivals when it split; it revealed
itself as one object with several honest faces, and the agreement of those faces is the strongest evidence the
construction can offer. When the book finally holds its bracket open --- the interval floored at $\alpha$ --- the
reason to believe it is not that this volume argued carefully. It is that one residue, read on four decorrelated faces
in four disjoint languages, returns the same interval. The jar is one residue read many ways, and that is why it can be
trusted. The next chapter takes this single, proven object and does the one thing left to do with it: it gives it a
name.
